\chapter{اللقاء الخامس والثلاثون: الصلاة الوسطى وتتمة آيات المتاع وبداية قصص البقرة}

\section{مقدمة}
السلام عليكم ورحمة الله وبركاته. الحمد لله رب العالمين، وأشهد أن لا إله إلا الله وحده لا شريك له، وأشهد أن محمداً عبده ورسوله، صلى الله عليه وعلى آله وصحبه وسلم.

استأنف الشيخ هذا اللقاء من قوله تعالى: \begin{ayah}حَافِظُوا عَلَى الصَّلَوَاتِ وَالصَّلَاةِ الْوُسْطَى\end{ayah}، ثم انتقل بعد ذلك إلى تتمة آيات المتاع والطلاق، وختم المجلس بفتح باب القصص في قوله تعالى: \begin{ayah}أَلَمْ تَرَ إِلَى الَّذِينَ خَرَجُوا مِنْ دِيَارِهِمْ وَهُمْ أُلُوفٌ حَذَرَ الْمَوْتِ\end{ayah}.

ويمتاز هذا اللقاء بأنه جمع بين ثلاثة أبواب تبدو متباعدة في الظاهر، لكنها متصلة في الحقيقة: تعظيم شعائر الله، وحسن الوقوف عند حدوده في أحكام الأسرة، والاعتبار بالقصص التي تبين قدرة الله وحكمته وسننه في عباده.

\separator

\section{تأويل (حافظوا على الصلوات والصلاة الوسطى)}
\isnad{قال أبو جعفر:} أمر الله بالمحافظة على الصلوات كلها، ثم خص الصلاة الوسطى بمزيد تنبيه، ففسر \rawi{الطبري} المحافظة هنا بالمداومة عليها في أوقاتها، وتعاهدها، وترك السهو عنها حتى يخرج وقتها.

ونبّه الشيخ إلى أن هذا التفسير يرد العناية بالصلاة إلى أصلين:
\begin{itemize}
    \item حفظ الوقت.
    \item وحفظ القلب من التضييع والغفلة.
\end{itemize}

فليس المقصود مجرد الأداء الصوري، بل القيام بحق هذه الصلوات تعاهداً ومراقبة.

\begin{fawaid}[حفظ العبادة يبدأ بحفظ وقتها]
من أول معاني المحافظة على الصلاة أن تؤدى في وقتها، لأن إخراجها عن وقتها أو التساهل فيه أصل من أصول التضييع، وإن بقيت صورة العمل.
\end{fawaid}

\section{الخلاف في تعيين الصلاة الوسطى}
استوعب \rawi{الطبري} الأقوال المشهورة في تعيين الصلاة الوسطى:
\begin{itemize}
    \item فقيل: هي صلاة العصر.
    \item وقيل: هي صلاة الظهر.
    \item وقيل: هي صلاة المغرب.
    \item وقيل: هي صلاة الفجر.
\end{itemize}

ثم عرض حجج كل فريق، فذكر ما ورد في العصر من الأخبار الكثيرة، ومنها حديث الأحزاب: \begin{ayah}شَغَلُونَا عَنِ الصَّلَاةِ الْوُسْطَى\end{ayah}، وما ورد في الظهر من مراعاة شدتها على الناس في الهاجرة، وما قيل في المغرب من كونها متوسطة بين صلاتين لا تقصر ولا تزيد، وما قيل في الفجر من جهة القنوت.

ومن أجمل ما في هذا الموضع أن \rawi{الطبري} لا يكتفي بنسبة القول إلى قائله، بل يشرح علة القول ووجهه، حتى يفهم القارئ لماذا قال كل فريق بما قال.

\begin{fawaid}[من تمام الانتفاع بالخلاف فهم علله لا مجرد حفظ نتائجه]
لا يكفي أن يعلم الطالب أن في المسألة أربعة أقوال مثلاً، بل لا بد أن يفهم وجه كل قول وحجته. وهذه من خصائص قراءة الطبري: أنه يفسر الأقوال كما يفسر الآيات.
\end{fawaid}

\section{ترجيح الطبري وأثر الأخبار في الباب}
يظهر من كثرة ما ساقه \rawi{الطبري} من الآثار الصحيحة والمشهورة أن قول العصر هو الأقوى في هذا الموضع، لما جاء فيه من نصوص صريحة، ولأنها جُمعت لها شواهد متعددة من جهة السنة والآثار.

وفي هذا تدريب لطالب العلم على أن الترجيح لا يبنى على الاحتمال الذهني المجرد، بل على كثرة الشواهد، وصراحة الدليل، ومجموع ما ورد في الباب.

\separator

\section{تأويل (وقوموا لله قانتين)}
أشار الشيخ إلى أن بعض الأقوال ربطت بين تعيين الصلاة الوسطى وبين معنى القنوت بعدها، كما في من قال إنها الفجر لاقترانها بالقنوت عند طائفة. وهذا يبين أن بعض المفسرين يبنون الترجيح على اقتران الألفاظ والسياق، لا على الرواية المجردة فقط.

لكن بقاء هذه المناسبات لا ينسخ ما جاءت به الأخبار الصريحة، بل تكون من جملة ما يذكر في توجيه الأقوال وفهم مأخذها.

\begin{fawaid}[قد يستفاد من اقتران الألفاظ وجه للقول لا يلزم منه ترجيحه]
من المهم أن يفرق الطالب بين \textit{وجه القول} و\textit{قوة القول}. فقد يكون للقول مأخذ حسن من السياق أو الاقتران، لكن يبقى الدليل المرفوع الصريح أرجح منه.
\end{fawaid}

\separator

\section{تتمة آيات المتاع: (وللمطلقات متاع بالمعروف حقا على المتقين)}
عاد الشيخ بعد ذلك إلى ختام آيات المتاع، فذكر أن \rawi{الطبري} رجح أن قوله تعالى: \begin{ayah}وَلِلْمُطَلَّقَاتِ مَتَاعٌ بِالْمَعْرُوفِ\end{ayah} يتناول جميع المطلقات من حيث الأصل، وأن هذه الآية جاءت بياناً عاماً بعد أن وردت في الآيات السابقة صور مخصوصة للمطلقات قبل المسيس أو قبله مع تسمية الصداق أو عدمها.

فهذا من بديع الجمع في القرآن: يذكر الخاص في مواضع، ثم يأتي العام الذي يشمل الباب من جهة أخرى، فلا يعارض بعضه بعضاً، بل يفسر بعضه بعضاً.

ورد \rawi{الطبري} بذلك على من حمل الآية على سبب خاص أو على معنى ضيق، ورأى أن الأظهر أن الله أراد بها تقرير أصل المتاع للمطلقات، مع مراعاة ما سبق من تفصيل الصور.

\begin{fawaid}[العام المتأخر قد يأتي لجمع الصور التي تقدمت مفصلة]
من فقه القرآن أن يذكر بعض الأحكام مفصلة في صور معينة، ثم يعقبها بلفظ عام يجمع الباب ويؤكد أصله. وهذا لا يوقع تعارضاً، بل يزيد الباب إحكاماً.
\end{fawaid}

\section{أثر أسماء الله وأوصاف العباد في آيات الأحكام}
وقف الشيخ عند مثل قوله تعالى: \begin{ayah}حَقًّا عَلَى الْمُحْسِنِينَ\end{ayah} و\begin{ayah}حَقًّا عَلَى الْمُتَّقِينَ\end{ayah}، وبيّن أن هذا ليس تزييناً لفظياً، بل من أعظم ما ينبغي أن يعتني به الفقيه والمربي. فربط الحكم بالإحسان والتقوى يربي العبد على أن القيام بحقوق الله والخلق لا ينفصل عن مقام القلب.

وكذلك في ختام الآيات بقوله: \begin{ayah}كَذَلِكَ يُبَيِّنُ اللَّهُ لَكُمْ آيَاتِهِ لَعَلَّكُمْ تَعْقِلُونَ\end{ayah}، فإن الله لا يأمر بتقواه أمراً مبهماً، بل يبين حدوده أولاً، حتى تكون التقوى عن علم، والعمل عن بصيرة.

\begin{fawaid}[التقوى لا تنفصل عن العلم بحدود الله]
ليس المتقي هو الذي يكثر التشدد على نفسه مجرداً، بل المتقي حقاً هو الذي يعرف حدود الله أولاً، ثم يحمل نفسه على الوقوف عندها. ولهذا كان البيان من الله مقدمة لازمة للتقوى.
\end{fawaid}

\separator

\section{بداية القصص: (ألم تر إلى الذين خرجوا من ديارهم وهم ألوف حذر الموت)}
افتتح \rawi{الطبري} بعد ذلك مقطعاً جديداً من قصص البقرة، فبيّن أن قوله: \textit{ألم تر} هنا من رؤية القلب، أي: ألم تعلم. لأن النبي صلى الله عليه وسلم لم يشهد هؤلاء بعينه، لكن الله أعلمه بخبرهم.

ثم ذكر أن القدر المحكم الثابت في القصة هو:
\begin{itemize}
    \item أنهم خرجوا من ديارهم.
    \item وأنهم كانوا ألوفاً.
    \item وأن خروجهم كان حذراً من الموت.
    \item وأن الله أماتهم ثم أحياهم.
\end{itemize}

أما تفاصيل العدد والموضع وسبب الخروج على وجه التحديد، فورد فيها كثير من الروايات، وأكثرها من الأخبار الإسرائيلية أو مما نقل عن بعض السلف في تفسير القصة.

\section{فائدة منهجية: كيف نتعامل مع روايات بني إسرائيل؟}
أفاد الشيخ هنا فائدة منهجية نفيسة: أن روايات بني إسرائيل ثلاثة أنواع:
\begin{itemize}
    \item ما شهد القرآن أو السنة الصحيحة بصدقه، فيقبل.
    \item وما شهد القرآن أو السنة بكذبه، فيرد.
    \item وما سكت عنه الشرع، فلا يصدق ولا يكذب على وجه الجزم، ويجوز حكايته إذا لم يترتب عليه محذور.
\end{itemize}

وبيّن أن العبرة الأصلية في القصص القرآني قد تتم من غير الوقوف على تلك التفاصيل، لكن قد يذكر بعضها من جهة الاستئناس أو تتميم الصورة، لا من جهة بناء حكم شرعي عليها.

\begin{fawaid}[العبرة في القصص قد تتم من غير تفاصيل الروايات الإسرائيلية]
إذا ثبت أصل القصة في القرآن، تمت العبرة به. أما التفاصيل التي لم تثبت، فقد يذكرها المفسر استئناساً، لكن لا يجعلها ركناً في فهم العبرة، ولا يبني عليها حكماً.
\end{fawaid}

\section{العطاء والمنع كلاهما ابتلاء}
ربط الشيخ بين ختام هذا الموضع وقوله تعالى بعده: \begin{ayah}مَنْ ذَا الَّذِي يُقْرِضُ اللَّهَ قَرْضًا حَسَنًا\end{ayah} و\begin{ayah}وَاللَّهُ يَقْبِضُ وَيَبْسُطُ\end{ayah}، فقرر أن الله يبتلي بالعطاء كما يبتلي بالمنع:
\begin{itemize}
    \item فالعطاء ابتلاء بالشكر، والإنفاق، وحسن التوظيف.
    \item والمنع ابتلاء بالصبر، والأخذ بالأسباب، وحسن الظن بالله.
\end{itemize}

وهذه فائدة تربوية كبيرة في باب العلم والمال والقدرة وسائر المواهب؛ إذ ليس كل معطى مكرماً مطلقاً، ولا كل ممنوع مهاناً مطلقاً، بل كل واحد منهما داخل في الامتحان الإلهي.

\begin{fawaid}[من فقه الحياة أن ينظر إلى العطاء والمنع على أنهما امتحان]
لا يصح أن يقاس فضل العبد بمجرد ما أوتي أو ما منع، بل بالنظر إلى كيف عمل في حال البسط، وكيف صبر واتقى في حال القبض. فالعبرة بالفعل في الحالتين، لا بمجرد الصورة الظاهرة.
\end{fawaid}

\separator

خلاصة هذا اللقاء أن \rawi{الطبري} جمع فيه بين تحرير معنى المحافظة على الصلاة، وبيان قوة الأخبار في تعيين الصلاة الوسطى، ثم رد آيات المتاع إلى أصول الإحسان والتقوى والعقل، وافتتح باب القصص القرآني بما يعلم القارئ كيف يميز بين المحكم من القصة، وما يذكر استئناساً من تفاصيلها.
